\section{CMB: a conditional compressed-geometry diagnostic}\label{sec:cmb}

\subsection{Data: the compressed CMB likelihood}
We use the compressed Planck~2018\,+\,ACT~DR6 likelihood of \citet{Bansal2025}, a three-parameter data vector: $R=100\sqrt{\Omega_bh^2+\Omega_{\rm cdm}h^2+\Omega_{\nu,\rm m}h^2}\;D_M(z_{\rm rec})/c$ (their Eq.~4), $\ell_a=\pi D_M(z_{\rm rec})/r_s(z_{\rm rec})$ and $\omega_b$, with $(R,\ell_a,\omega_b)=(1.7504,\,301.77,\,0.022371)$ and the published $3\times3$ covariance. We retain the declared Bansal--Huterer vector rather than replacing it with the newer CMBComp product \citep{CMBComp2026}: the latter uses a different SPA dataset and compression protocol, remains GR-based and therefore does not remove the modified-gravity transplant caveat. It is not a drop-in replacement for the present input; a separate robustness comparison is deferred.
Massive neutrinos ($\sum m_\nu=0.06$~eV) are treated with a two-asymptote scheme: non-relativistic matter, $\omega_\nu(1+z)^3$, below the join, and the relativistic density inside $N_{\rm eff}=3.046$ of the radiation term above it, with the two forms matched continuously at their crossing, $1+z=\omega_\nu/\omega_{r,\nu}=37.7$ ($h$-independent; $kT_\nu\simeq6.3$~meV there, i.e.\ $T\sim m_{\nu,i}/3$ for an individual mass $m_{\nu,i}\sim0.02$~eV). The residual transition-region approximation affects $E^2(z)$ at the $10^{-4}$ level absolute (in $E^2$ itself, not relative to any density fraction), below every scale discussed here.
The recombination redshift $z_{\rm rec}$ is evaluated from the \citet{HuSugiyama1996} fitting formula at the fitted $(\omega_b,\omega_m,h)$ rather than held fixed at $1090$. Although the published mean vector and covariance were initially extracted at fixed $z_*=1090$, \citet[][Appendix~A]{Bansal2025} also validate the alternative with CAMB-determined $z_*$. Our fitted-$z_{\rm rec}$ choice follows that alternative under the assumption of standard recombination physics; the Hu--Sugiyama formula is not applied here to an exotic thermal history. A separate fixed-$z_*=1090$ sensitivity run is not reported.
\paragraph{Applicability to modified gravity.}
The compressed likelihood is derived under GR, and \citet{Planck2015XIV} warn that it may be inapplicable where perturbation growth differs from GR. We use it only conditionally. Vainshtein screening suppresses the background modification at $z_{\rm rec}$, where $H(z_{\rm rec})\gg\mFP$, but screening alone does not validate the transfer of a GR-calibrated likelihood to bimetric perturbations. The negativity of the background diagnostic $X$ at $z>z_*$ (\cref{sec:scalar}) is a separate statement from the applicability of linear perturbation theory: screening does not make the background $X$ nonnegative, but its effect on the actual evolution of the linear instability is not established here. Our CMB results are therefore a background-level geometric diagnostic with this caveat. Recent direct full-versus-compressed comparisons show that compression can fail when non-standard dark-matter evolution changes the density mapping relevant to the CMB \citep{Pedrotti2026}; the true/effective-density split studied here is precisely why we do not promote the compressed result to a full bimetric CMB constraint.

\subsection{Results}\label{sec:cmbres}
The comparative reference is a flat \LCDM{} model with three free physical parameters $(\omega_m,h,\omega_b)$; $\chi^2_{\rm ref}=0$ to numerical precision at $\omega_m=0.14298$, $h=0.67568$, $\omega_b=0.022371$ (the reference carries the massive-neutrino density in the $R$ prefactor and in the $z_{\rm rec}$ argument, as does the bimetric path). No reference goodness-of-fit $p$-value is therefore defined, and the comparative $\Dchi$ equals the bimetric absolute $\chi^2$ (DOF\,=\,1).
Within this conditional geometric test, rather than a full bimetric CMB likelihood, the bimetric baseline gives $\chi^2=0.0001703$ at $h_*=0.6869$ ($p=0.990$), consistent with the compressed data and, notably, with the $H_0=68.84\ \mathrm{km\,s^{-1}\,Mpc^{-1}}$ inferred from the BAO distance fit (\cref{sec:Om-H0}). The two determinations use disjoint distance data, but they are not fully independent: the BAO inference adopts the Planck (GR-calibrated) sound horizon $r_d=147.09$~Mpc, so the agreement is a consistency check of the background fit, not an independent confirmation. Neither the absolute $p$-value nor the small $\Dchi$ constitutes a model-selection significance. Thus the benchmark itself carries no appreciable CMB penalty; the scan dependence below is conditional on the adopted transplant protocol (\cref{tab:scans} and \cref{tab:A5}):
\begin{itemize}
\item \emph{$(\alpha,\beta)$ scan:} $\Dchi=0.0000$--$0.0053$ across the 19 non-degenerate combinations; the compressed-CMB baseline is essentially blind to the potential parameters at $\theta=6.1\deg$.
\item \emph{$\mFP/H_0$ scan} ($5$ to $100$): $\Dchi=0.0002$--$0.2170$.
\item \emph{$\theta$ scan:} $\Dchi$ rises from $0.0002$ at $6.1\deg$ to $2412.98$ at $30\deg$. At $30\deg$ with $\mFP/H_0=25$ the late-time de~Sitter form of the Higuchi bound is satisfied ($\mFP^2/H_0^2=625$ against $2\OL=1.4$ at the benchmark anchor and $1.68$ at the $30\deg$ minimum), although the angle is strongly penalised within the adopted $\lambda=0$ compressed-CMB geometric diagnostic. This $\theta$-sensitivity is conditional on the anchor and protocol: \cref{sec:anchor} shows that at fixed $\theta$ the anchor is tightly configured and that freeing $\OL$ does not remove the joint tension but relocates it.
\end{itemize}

\paragraph{Split test.}
The split test isolates the source of the $\theta$ penalty in the $D_M$ integral. The bimetric residual relative to the screened background, $E^2_{\rm bim}(z)-E^2_{\rm scr}(z)$, is removed above $z=12$ (\emph{zero\_high}) or below $z=12$ (\emph{zero\_low}), with the screened sound horizon retained. The diagnostic is reported through the shift parameter $R$ only. The published correlations are $\rho(R,\omega_b)=-0.637$, $\rho(\ell_a,\omega_b)=+0.220$ and $\rho(R,\ell_a)=-0.040$. Moreover, the two modes modify a nonlinear distance integral while sharing the full-mode $\ell_a$, so they are not independent likelihood components and their mode $\chi^2$ values are neither additive nor refit goodness-of-fit measures. Those legacy mixed-slot fields remain in the saved output for audit, but are not tabulated or interpreted. At $30\deg$, $R$ moves from $1.94441$ (full) to $1.93500$ (zero\_high) and $1.78743$ (zero\_low), against the observed $1.7504$: the late-time unscreened region supplies the dominant share of the geometric mismatch, and the deep-past negative asymptote of $\ODE$ contributes a few percent. At $10\deg$ the same pattern holds at small amplitude ($R=1.76330$ full, $1.76210$ zero\_high, $1.74789$ zero\_low).

\paragraph{The transplant family $R(\lambda,\theta)$.}
The $R$ prefactor contains the true matter density, whereas the late-time background responds to the effective combination $\cos^2\!\theta\,\Om$ (the dark degeneracy, \citealp{Kunz2009}); in GR the two coincide, and the compressed likelihood was derived under GR. This transplant is a known hazard: the model dependence of the shift parameter was flagged by \citet{Elgaroy2007}; \citet{Planck2015XIV} exclude modified-gravity models from their compressed likelihood precisely because of this degeneracy; and \citet[][Appendix~B]{Hogas2021b} noted the same danger within bimetric gravity itself: in the GR limit $R$ carries a $1/\cos\theta$ factor and ``will be far away from the tabulated value even if the background expansion follows the best-fit \LCDM{} model''. What we add beyond these warnings is the one-parameter family $R(\lambda,\theta)$, the explicit $\lambda^*(\theta)$ sweep and the structural saturation argument with grid-supported estimates. The interpolating family considered here has one parameter. A clarifying word on
the status of $\lambda$ is in order: $\lambda$ is a parameter of this
diagnostic transplant of a GR-compressed likelihood onto a dark-degenerate
background, \emph{not} a parameter of the Hassan--Rosen theory; the saturation
below is a property of this particular three-data/three-parameter mapping, not
evidence of a physical equivalence of theories, and a minimum over $\lambda$
is not a measurement of a new physical quantity. Which density physically enters $R$ is a genuine physical question; the adopted filling is a protocol, not an answer, and the GR-calibrated compression cannot settle it. The three readings of the family (the true-density map $\lambda=0$, the effective map $\lambda=1$, and the profiled vertex $\lambda^*$) are diagnostic protocols, not three physical variants of the theory. The family reads:
\begin{equation}
\begin{aligned}
R(\lambda,\theta)&=100\,h\sqrt{(1-\lambda\sin^2\!\theta)\,\Om+\omega_\nu/h^2}\;\frac{D_M(z_{\rm rec};E^2_{\rm bim})}{c}\\
&=100\sqrt{(1-\lambda\sin^2\!\theta)\,\omega_m+\omega_\nu}\;\frac{D_M(z_{\rm rec};E^2_{\rm bim})}{c},\qquad
\omega_m=\Om h^2,\quad\lambda\in[0,1],
\end{aligned}\label{eq:Rlambda}
\end{equation}
with $\lambda=0$ the quartic (true-density) map and $\lambda=1$ the effective map. $D_M$ does not depend on $\lambda$. The data fix the combination $(1-\lambda\sin^2\!\theta)\Om$, so at each $\theta$ they are satisfied wherever
\begin{equation}
\lambda^*(\theta)=\bigl[1-\Om^{\rm data}(\theta)/\Om^{\rm quartic}(\theta)\bigr]/\sin^2\!\theta
\end{equation}
falls inside $[0,1]$. Here $\Om^{\rm data}(\theta)$ is the effective matter density preferred by the compressed data at angle $\theta$, operationally the density at which the fixed-$\lambda$ penalty vanishes, computed by the sweep (\code{results/cmb\_lambda\_theta\_sweep.json}).
At the benchmark angle the parabolic estimate $\lambda^*$ is negative: the data prefer a density above the quartic one. The five-point sweep shows $\lambda^*$ crossing zero between $6.1$ and $10\deg$; local three-point parabolas estimate $\lambda^*=0.48$ at $10\deg$ and $0.75$ at $30\deg$ (\cref{tab:A5}, \cref{fig:cmbtransplant}). Wherever $\lambda^*$ is admissible, the three data points can be exhausted by the three parameters $(h,\omega_b,\lambda)$, so saturation is structural in this mapping. The quoted vertex locations and residuals are nevertheless grid-limited numerical estimates, not a continuous optimisation certificate: the parabolic residual at $10\deg$ is $0.0086$, while the best of the five sampled $\lambda$ values remains below $\sim2.5$ for $\theta\ge10\deg$. Numerically ($\chi^2$ at $6.1\deg$ / $30\deg$): $\lambda=0$ gives $0.0001703$ / $2412.98$; $\lambda=0.75$ at $30\deg$ gives $0.56$; $\lambda=1$ gives $6.2871$ / $275.61$.
The apparent $\theta$-sensitivity of any fixed map is the drift of $\lambda^*(\theta)$. Within the scanned fixed-map family, the apparent $\theta$ constraints are therefore unstable to the transplant choice; the remaining model-dependent diagnostic is the effective-law residual
\begin{equation}
\delta E^2(z)=E^2_{\rm bim}-\bigl[\OL+\cos^2\!\theta\,\Om(1+z)^3+\Orad(1+z)^4\bigr]=\sin^2\!\theta\,\mFP^2\bigl[\sin^2\!\theta\,u-G(\delta)\bigr]\quad\text{(exact)},
\end{equation}
which is $1.58\times10^{-5}$ of $E^2$ at $z=5$ and $3.6\times10^{-3}$ at $z=12$ for the BAO closure used in \cref{fig:mechanism}a. The SNe/scalar $N_{\rm eff}$ closure gives $1.43\times10^{-5}$ at $z=5$; this convention shift does not affect any conclusion. The near-cancellation between the two terms is the analytic content of the dark degeneracy, not an accident.

\begin{figure}[t]
\centering
\includegraphics[width=\textwidth]{fig5_cmb_transplant}
\caption{The $R(\lambda,\theta)$ transplant of the compressed CMB likelihood. (a)~$\chi^2$ against $\lambda$ at five angles (numbers label the curves; the sweep samples five values of $\lambda$, so the sharp minima are the grid resolution). (b)~Parabolic vertex $\lambda^*(\theta)$ through the three sweep points around the minimum; it is negative at $6.1\deg$ (open marker, boundary) and admissible ($0\le\lambda^*\le1$, shaded) from $10\deg$. (c)~$\chi^2$ against $\theta$ for the true-density map ($\lambda=0$), the effective map ($\lambda=1$) and the best of the five sampled values of $\lambda$ (below $\sim2.5$ for $\theta\ge10\deg$): the apparent $\theta$-penalty of any fixed map is the drift of $\lambda^*(\theta)$. Generated by \code{src/make\_figures.py} from \code{results/cmb\_lambda\_theta\_sweep.json}.}\label{fig:cmbtransplant}
\end{figure}

\paragraph{Adopted protocol.}
For dark-degenerate backgrounds the $R$ slot is filled with the true density ($\lambda=0$): the physically motivated endpoint, since the equality-scale physics responds to the true density, and the convention against which the $\lambda$-systematic is quoted (the sound horizon is computed with the true density). A principled re-derivation of the compression for this benchmark awaits a Boltzmann-level treatment. Partial bimetric perturbation calculations exist (tensor modes have been evolved and found to grow in parts of the parameter space, \citealp{Johnson2015,Amendola2015}), but they neither validate the GR-calibrated compression nor substitute for it; a standard early universe was likewise assumed as a working premise in earlier bimetric data analyses \citep{Caravano2021}. \Cref{sec:anchor} shows that no stronger background-level statement is possible within the scanned plane.

\subsection{Summary of scan results}\label{sec:scansummary}
\Cref{tab:scans} collects the per-probe scans of \cref{sec:bao,sec:sne,sec:cmb} on one page: each row varies one ingredient of the benchmark ($(\alpha,\beta)$, $\mFP/H_0$ or $\theta$) and reports the resulting range of $\Dchi$ for each probe. BAO and SNe are flat across the potential parameters and the mass scale. In the self-consistent fixed-anchor $\theta$ scan, BAO and SNe are nearly flat, whereas the large CMB variation is conditional on the adopted $R$-slot convention (\cref{sec:cmbres}). The verdicts in the last row are the per-probe conclusions used in \cref{sec:anchor,sec:conclusions}.

\begin{table}[H]
\centering\small
\caption{Summary of scan results (per-probe $\Dchi$; CMB against the free three-parameter \LCDM{} reference, i.e.\ the absolute bimetric $\chi^2$).}\label{tab:scans}
\begin{tabularx}{\textwidth}{@{}l>{\raggedright\arraybackslash}X>{\raggedright\arraybackslash}X>{\raggedright\arraybackslash}X@{}}
\toprule
Scan & BAO & SNe & CMB\\
\midrule
$(\alpha,\beta)$, 19 points & $0.1791$--$0.1794$ & $7.315$ (diag., flat), $2.936$ (full, flat) & $0.0000$--$0.0053$\\
$\mFP/H_0$ & $0.179$ (flat) & $7.315$ (diag., flat), $2.936$ (full, flat) & $0.0002$--$0.2170$\\
$\theta$ (baseline $6.1\deg$) & $0.179$ (self-consistent) & $7.315$ (diag.); $2.936$ (full) & $0.0002$\\
$\theta$ (range $6.1$--$31.5\deg$; CMB to $30\deg$) & $0.179\to0.177$ (self-consistent); fixed-$\Om$ variant reaches $55.3$ & $7.315\to7.320$ (diag.); $2.936\to2.939$ (full) & $\to2412.98$ at $30\deg$\\
Verdict & Blind & Per-probe offset (parameter mismatch) & No robust $\theta$ information (transplant convention)\\
\bottomrule
\end{tabularx}
\end{table}
